% 原文 p.65（後半）


\subsubsectionlink{経路を監視するものとしての環境}
この本を考えよう。本は、空気分子、可視光や熱の光子、さらには宇宙のニュートリノや放射性の背景放射からなる、大きな環境に浸っている。毎秒、膨大な数のこれらの粒子が本に衝突し、散乱される。各衝突は、本の位置と向きに応じて、粒子をある程度偏向させる。ある一個の空気分子に注目しよう。それは、ある初期位置から一定の速度で出発し、本に散乱され、入射経路に対して一定の角度だけ曲がった軌道に沿って飛び去る。

本の向きが違う角度だったとしよう。その場合、空気分子は一般に、最初の場合とは異なる方向へ飛び去る。したがって、散乱された分子のこの二つの異なる経路によって、本に二つの異なる空間的向きが存在することを区別できる。言い換えれば、散乱粒子である空気分子は、ここでは本である散乱物体の位置と向きについての情報を持ち去る。

% 原文 p.66
\hypertarget{srcpage:66}{}\label{srcpage:66}
すでになじみ深い用語を使えば、それらは系についての経路情報を符号化する。ここで、以後の本文でも同じだが、「経路情報」という語は、より一般的に「どの状態かという情報」の意味に理解してほしい。前者の用語は、関心のある情報がスリットを通る粒子の経路に関わる、二重スリット実験の例によって、歴史的に動機付けられている。それでも、実際には粒子の軌道に関する情報でない場合にも、定着した用語「経路情報」をしばしば使う。

本に衝突する一個の空気分子は、本の向きについて、多くの情報を持ち去らないかもしれない。結局、散乱が起こるのは、本の大きさに比べて微小な領域だけである。しかし上に述べたように、どの瞬間にも、何百万個もの分子が本の表面に散乱されている。これらすべての分子を合わせたものに符号化される経路情報、あるいは、より適切には「どの向きかという情報」は、本の二つの異なる位置と向きがとても似ている場合でも、それらを識別するのに、確かに十分であろう。

この議論を、量子形式によって模式的に表そう。1秒間に本で散乱される$N$個すべての粒子に注目する。衝突前の$N$個すべての環境粒子の量子状態を$\lvert E_0\rangle$で表す。衝突後、各粒子は何らかの速度で、ある軌道に沿って飛び去る。本の向きが別であれば、粒子は少し異なる方向へ散乱されるかもしれない。$N$個すべての粒子が、本に散乱されたとする。本の向きに応じて、散乱後の環境状態を、それぞれ$\lvert E_1\rangle$、$\lvert E_2\rangle$と書く。

上の議論によれば、この二つの状態に含まれる情報は、本の二つの向きを識別するのに十分なので、これらは明らかに区別できる。したがって、$\lvert E_1\rangle$と$\lvert E_2\rangle$の重なりは、無視できるほど小さい\footnote{$\lvert E_1\rangle$、$\lvert E_2\rangle$を、個々の分子に属する状態の積、$\lvert E_1\rangle=\prod_{i=1}^N\lvert e_1\rangle_i$、$\lvert E_2\rangle=\prod_{i=1}^N\lvert e_2\rangle_i$と考えられる。各分子状態の対の重なりは1に近く、$(\langle e_1\vert e_2\rangle)_i\approx1$であるが、そのような個別状態を多数$N$個掛け合わせた積状態$\lvert E_1\rangle$、$\lvert E_2\rangle$は、近似的に直交する。}。前の2.6.2節で導入した経路情報の形式とのつながりが、明瞭に分かる。環境状態$\lvert E_1\rangle$、$\lvert E_2\rangle$は、単に経路検出器、ここでは向き検出器の状態に対応する。

重要なのは、このような環境の経路検出器が、自然界のあらゆる場所にあることである。どの物体も環境と相互作用し、環境は系のある物理的性質について情報を得る。本を高い真空度の中に置けば、空気分子をある程度遮蔽できるかもしれない。可視光の光子を遮断し、本を冷やして熱光子の影響も減らそうとできるかもしれない。それでも、次の第3章で見るように、他の源からの粒子が残り、系の位置を「測定」し続ける。系に対して、von Neumannの意味での実効的測定を連続的に行う、遍在する「測定装置」としての環境の、この決定的役割を、初めて明瞭に認識して論じたのは、1980年代初頭のZurekであった\cite{ref008,ref076}。

% 原文 p.67
\hypertarget{srcpage:67}{}\label{srcpage:67}
この監視過程に、いかなる人間の観測者も必要ないことを強調する。環境に符号化された経路情報を、原理上読み出せるという事実だけで、干渉模様が消え、粒子が「波としての性質を失う」のに十分である。これにより、2.1.3節ですでに引用した、Heisenbergの「粒子の軌道は、私たちがそれを観測するという行為によって作られる」\cite[p.~185]{ref042}という言葉に、より正確で、観測者に依存しない意味が与えられる。この実証主義的態度は、第1章で述べたEinsteinの有名な修辞的な問い、「月は私が見ているときにしか存在しないのか」\cite{ref003}に代表されるように、量子論に対する多くの批判を生んできた。

ここで、2.6.2節と同じ形式的議論を行える。系が、異なる二つの位置$x_1$、$x_2$のまわりへの局在に対応する量子状態$\lvert\psi_1\rangle$、$\lvert\psi_2\rangle$の、コヒーレントな重ね合わせで記述されるとしよう。二重スリット実験では、対応する位置空間の波動関数$\psi_1(x)$、$\psi_2(x)$は、スリットでの部分波を表す。環境粒子の散乱が起こる前、系と環境を合わせた状態は、積の形
\begin{equation}\lvert\psi\rangle\lvert E_0\rangle=\frac1{\sqrt2}(\lvert\psi_1\rangle+\lvert\psi_2\rangle)\lvert E_0\rangle\tag{2.70}\end{equation}
を持つ。上の議論に従い、系と環境の相互作用の力学は、
\begin{align}\lvert\psi_1\rangle\lvert E_0\rangle&\longrightarrow\lvert\psi_1\rangle\lvert E_1\rangle,\tag{2.71a}\\\lvert\psi_2\rangle\lvert E_0\rangle&\longrightarrow\lvert\psi_2\rangle\lvert E_2\rangle\tag{2.71b}\end{align}
で与えられる。すなわち、環境の状態は、系の状態に応じて$\lvert E_1\rangle$または$\lvert E_2\rangle$に発展する。するとSchrödinger方程式の線形性から、通常のvon Neumann発展
\begin{equation}\frac1{\sqrt2}(\lvert\psi_1\rangle+\lvert\psi_2\rangle)\lvert E_0\rangle\longrightarrow\frac1{\sqrt2}(\lvert\psi_1\rangle\lvert E_1\rangle+\lvert\psi_2\rangle\lvert E_2\rangle)\tag{2.72}\end{equation}
が従う。系の相対状態$\lvert\psi_1\rangle$、$\lvert\psi_2\rangle$は、経路情報を符号化する環境状態$\lvert E_1\rangle$、$\lvert E_2\rangle$ともつれた。初めは系だけに限られていた重ね合わせが、より大きな、系と環境の複合状態へ広がった。それに対応して、成分$\lvert\psi_1\rangle$、$\lvert\psi_2\rangle$間のコヒーレンスは、もはや系だけの性質ではなく、大域的な系と環境の状態の共有された性質となった。そのため、コヒーレンスが、環境を含むようになった「より大きな系の中へ非局在化した」\cite[p.~5]{ref077}と言うことが多い。式(2.72)で記述される、系と環境の力学的発展が、デコヒーレンス過程の基本的な形式的表現であり、その帰結を次に論じる。

% 原文 p.68
\hypertarget{srcpage:68}{}\label{srcpage:68}
\subsectionlink{デコヒーレンスと干渉の局所的減衰}
状態(2.72)における系の縮約密度行列$\hat\rho_{\mathrm{particle}}$は、式(2.49)、(2.66)より
\begin{equation}\begin{split}\hat\rho_{\mathrm{particle}}=\frac12\{&\lvert\psi_1\rangle\langle\psi_1\rvert+\lvert\psi_2\rangle\langle\psi_2\rvert\\&+\lvert\psi_1\rangle\langle\psi_2\rvert\langle E_2\vert E_1\rangle+\lvert\psi_2\rangle\langle\psi_1\rvert\langle E_1\vert E_2\rangle\}\end{split}\tag{2.73}\end{equation}
で与えられる。いつものように、最後の二項は成分状態$\lvert\psi_1\rangle$、$\lvert\psi_2\rangle$の間の干渉に対応する。環境が実際に十分な経路情報を記録したなら、最終的な環境状態$\lvert E_1\rangle$、$\lvert E_2\rangle$は近似的に直交し、$\langle E_2\vert E_1\rangle\approx0$になる。上の、空気分子が巨視的物体に1秒間にわたって散乱される例では、これは確かに成り立つ。そのとき、縮約密度行列(2.73)の干渉は抑制される。
\begin{equation}\hat\rho_{\mathrm{particle}}\approx\frac12\{\lvert\psi_1\rangle\langle\psi_1\rvert+\lvert\psi_2\rangle\langle\psi_2\rvert\}.\tag{2.74}\end{equation}
系と環境の両方を含む測定だけが、重ね合わせ状態(2.72)の成分間に存続するコヒーレンスを、明らかにしうる。しかし実際には、系と相互作用した多数の環境自由度のすべてを観測に含めることは不可能である。JoosとZeh\cite[p.~224]{ref007}は、この状況を「干渉項はなお存在するが、そこにはない」と端的にまとめた。すなわち、干渉項は、系と環境の重ね合わせ(2.72)という大域的水準には残るが、縮約密度行列(2.74)で記述される系の局所的水準では、観測不能となる。もちろん通常、環境自由度の一部は、観測の一部である。たとえば系から散乱された光の一定部分を、直接捕らえる\footnote{実際、通常、観測者が系の情報を集めるのはこのような方法による（2.9節）。}。その場合、これらの環境自由度を系の一部と見なせばよい。しかし重要なのは、直接には観測しない、他の環境自由度が、なお非常に多数残ることである。

系と相互作用する環境自由度が非常に多く、私たちにはそれを直接操作できないため、式(2.72)で記述される系と環境のもつれの生成を、実際に元へ戻すことは、ほとんど不可能である。そのため環境に誘起された局所的な位相コヒーレンスの喪失、すなわち、干渉効果を観測するために必要な、重ね合わせの成分間の明確な位相関係の喪失は、通常、実際上のあらゆる目的に対して不可逆である（真に、または実効的に可逆なデコヒーレンスは2.13節で論じる）。

% 原文 p.69
\hypertarget{srcpage:69}{}\label{srcpage:69}
不可避で遍在する、環境による監視が引き起こす、系と環境の複合状態への位相関係の、実際上不可逆な非局在化こそが、デコヒーレンス過程である。環境が誘起する優先基底の選択は、このデコヒーレンス過程の帰結と見なすことができ、後の2.8節で論じる。この過程は、局所系に実効的に非ユニタリな力学をもたらし、それは、たとえば干渉模様の減衰として現れる［式(2.74)］。多数の部分系からなる環境は、系の状態に対する、増幅する高次の測定装置として働く。測定の定義自体が、不可逆性という性質に依拠することに注意しよう。

したがって環境による監視と、その結果のデコヒーレンス過程は、2.5.3節で論じた、干渉が観測されない問題への解決を与える。すでに論じたように、重ね合わせ状態の成分間の干渉効果を観測するためには、その成分が測定されていてはならず、どの状態かという情報が利用可能であってはならない。しかし、自然界に見いだすすべての物理系は、周囲と強く相互作用する開いた量子系である。周囲は連続的に系の情報を獲得し、コヒーレンスが系から環境へ絶えず「漏出」する。環境粒子の散乱の例では、系が大きいほど、単位時間当たりに多くの粒子が表面で跳ね返り、経路情報を持ち去る。したがって、位置空間で十分に離れた波束の重ね合わせで記述される、中間的・巨視的物体は、通常、きわめてデコヒーレンスを起こしやすい。そのことが、そのような重ね合わせと、対応する空間的干渉模様を、自然界で観測するのが、なぜそれほど難しいかを説明する。

2.4.6節（原文48頁）で述べた、縮約密度行列を状態の真正混合を記述するものと誤解しないように、という注意を繰り返す。縮約密度行列(2.74)は、混合状態の密度行列と形式上同じでも、非真正混合だけを記述する。したがって、式(2.74)から、系が実際に二状態$\lvert\psi_1\rangle$、$\lvert\psi_2\rangle$の一方にあると結論してはならない。結局、両成分は、複合量子状態(2.72)に、完全に等しく残っている。環境との純粋にユニタリな相互作用だけでは、いずれか一成分を選び出すことは、単にできない。この問題は8.1節でさらに詳しく分析する。

デコヒーレンスについての上の説明は、環境の相対状態$\lvert E_1\rangle$、$\lvert E_2\rangle$がどのように識別可能となり、環境が系の経路検出器として働くかについての、かなり直観的な議論に基づいた。デコヒーレンス研究の大部分は、さまざまな関心のある物理的状況について、系と環境の相互作用を現実的に模型化することに向けられる。それは環境状態の時間発展、および、その結果として系の縮約密度行列の力学に及ぼす影響について、精密な力学的記述を与える。通常、こうした模型は、環境が監視する異なる系状態$\lvert s_i\rangle$と量子的に相関した、異なる環境の相対状態$\lvert E_i(t)\rangle$が、系と環境の相互作用の結果、ごく短い時間尺度で識別可能になることを示す。

% 原文 p.70
\hypertarget{srcpage:70}{}\label{srcpage:70}
具体的には、後に見るように、多くの系と環境の模型で、異なる環境の相対状態$\lvert E_i(t)\rangle$の重なりは、指数減衰に従う。すなわち$i\ne j$に対し、
\begin{equation}\langle E_i(t)\vert E_j(t)\rangle\propto e^{-t/\tau_d}.\tag{2.75}\end{equation}
ここからは、相互作用を「入れる」時刻を$t=0$とする。$t<0$では通常、系と環境は完全に無相関と仮定する。量$\tau_d$は特徴的なデコヒーレンス時間尺度であり、各模型のパラメータを具体的に選べば、数値的に評価できる。

環境状態$\lvert E_i(t)\rangle$の重なりの、関数形、関係する時間尺度などの減衰特性は、式(2.73)を介して、これらと一対一に相関した系の重ね合わせ状態の成分$\lvert s_i\rangle$間のコヒーレンスの抑制について、その特徴的時間依存性と時間尺度に、そのまま反映される［式(2.72)］。したがって、与えられた系と環境の模型のデコヒーレンス力学を定量的に研究するには、$\lvert E_i(t)\rangle$の具体形を求めなければならない。後の2.10節で最初の例を調べ、より現実的な模型は本書の後半で扱う。

特に第3章では、空気分子や光子などの環境粒子の散乱が誘起する、空間的デコヒーレンスの具体的な模型化を行う。先取りすれば、環境の相対状態の重なりが
\begin{equation}\langle E_x(t)\vert E_{x'}(t)\rangle\propto e^{-\Lambda\lvert x-x'\rvert^2t}\tag{2.76}\end{equation}
のように減衰することが分かる。ここで$\lvert E_x(t)\rangle$は、系の重心が位置$x$に局在することを記述する系の状態と、一対一に相関した環境の相対状態である。定数$\Lambda$は、散乱過程の具体的な物理的性質で決まる。重なり(2.76)は、時間とともに指数関数的に減少する。したがって、式(2.76)は、しばしば現れる指数減衰の振る舞い(2.75)の一例である。環境粒子が物体を連続的に散乱し、環境に符号化される経路情報、すなわち系の位置に関する情報量を、着実に増加させる。

式(2.76)からは、特徴的なデコヒーレンス時間$\tau_d$［式(2.75)］が、系の異なる重心位置$x$、$x'$の間のコヒーレントな距離$\lvert x-x'\rvert$の二乗に、反比例することも分かる。この結果は、次の直観的な議論で理解できる。第3章では、より厳密に説明する。

% 原文 p.71
\hypertarget{srcpage:71}{}\label{srcpage:71}
位置$x$と$x'$が近ければ、散乱された環境粒子の、対応する最終的な相対状態は、とても似ていると予想でき、環境が二位置の空間的な違いを分解するのは難しい。したがって重なり$\langle E_x(t)\vert E_{x'}(t)\rangle$が顕著に減衰するには、各々が非常に不完全な経路情報だけを符号化する散乱事象を、多数蓄積する必要がある。逆に重心位置$x$、$x'$が離れるほど、環境はこの二位置を容易に識別できる。各散乱事象が、二つの重心位置の差を完全に分解できるほど$\lvert x-x'\rvert$が大きい極限では、重なり$\langle E_x(t)\vert E_{x'}(t)\rangle$の指数減衰率は、距離$\lvert x-x'\rvert$に依存しなくなり、全散乱率$\Gamma_{\mathrm{tot}}$に比例するだけになる。この極限では、式(2.76)は
\begin{equation}\langle E_x(t)\vert E_{x'}(t)\rangle\propto e^{-\Gamma_{\mathrm{tot}}t}\tag{2.77}\end{equation}
に置き換えられる。式(2.76)、(2.77)は第3章で導出して論じる。他の重要なデコヒーレンス模型と、その力学は、第5章で詳しく論じる。
